\section{Verification purpose}

\textbf{Case ID:} \CaseID

This case verifies sensitivity of firebrand-driven WUI spread to structural spatial representation. Five variants cover the baseline grid, native fine and coarse grids, and aligned and misaligned integrated-structure representations under one physical geometry and source definition.

The test quantifies native cell-scale distortion and determines whether integrated representation restores structure-level accumulation, ignition chronology, and spread-rate consistency. Capability-gated variants remain explicit and are not replaced by a physically different native-cell calculation.

\section{Mathematical formulation and numerical implementation}

\subsection{Heat release and generation}

For a structure ignited at \(t_i\), the design heat-release-rate per unit area is
\[
q''(\tau)=
\begin{cases}
q''_{\max}\tau/300,&0<\tau<300,\\
q''_{\max},&300\le\tau\le3900,\\
q''_{\max}(4200-\tau)/300,&3900<\tau<4200,\\
0,&\text{otherwise},
\end{cases}
\quad \tau=t-t_i,\qquad q''_{\max}=\SI{400}{kW.m^{-2}}.
\]
For footprint \(A_s\), \(\mathrm{HRR}_i=q''A_s\), and the emission rate is
\[
\dot N_i(t)=G_R'\mathrm{HRR}_i(t),\qquad
G_R'=\SI{10}{pcs.MW^{-1}.s^{-1}}.
\]
The configuration therefore uses per-MW generation; per-area generation would
not preserve heat-release-scaled production for this design fire.

\subsection{Himoto transport and analytical load}

The downwind distance is lognormal, with physical-space moments
\[
\mu_x/D=0.47(B^*)^{2/3},\qquad \sigma_x/D=0.88(B^*)^{1/3},
\]
\[
B^*=\frac{U}{\sqrt{gD}}
\left(\frac{\rho_p}{\rho_\infty}\right)^{-3/4}
\left(\frac{d_p}{D}\right)^{-3/4}
\left(\frac{\mathrm{HRR}}
{\rho_\infty c_pT_\infty g^{1/2}D^{5/2}}\right)^{1/2}.
\]
Here \(D=\SI{10}{m}\), \(\rho_p=\SI{100}{kg.m^{-3}}\), and
\(d_p=\SI{5}{mm}\), matching the reference formulation and current source. The postprocessor
independently evaluates the analytical deposited-load convolution
\[
N_k=\sum_{i<k}\int_0^T\dot N_i(t)
\left[\int_{x_k}^{x_{k+1}}
f(x-x_i;\mu_i(t),\sigma_i(t))\,dx\right]dt
\]
with a \SI{5}{s} quadrature interval. Deposited counts are divided by physical
footprint area, not a resolution-dependent raster-cell area.

\subsection{Physical ignition and expected chronology}

The configured ignition model uses
\(N_{\rm crit}=\SI{2928.9}{pcs.m^{-2}}\)~\cite{Qin2025}. At the structure
spanning \(20\le x\le\SI{30}{m}\), this load is reached near
\SI{2496.9}{s}. Adding the \SI{42.1}{s} small-fuel transition and
\SI{300}{s} development delay predicts full ignition near \SI{2858.8}{s}.
Targets beyond the roughly \SI{80}{m} peak-HRR spotting range require
successive structure ignition and consequently have nonlinear arrival times.

\subsection{Resolution signatures}

On the \SI{2.5}{m} native mesh (native fine-grid variant), one physical structure occupies
multiple independent cells. Uneven deposition and cell-local ignition should
produce a nonzero within-structure load range and mean ROS near
\SI{0.001}{m.s^{-1}}, compared with about \SI{0.006}{m.s^{-1}} for the
exact-grid baseline.
On the \SI{30}{m} native mesh (native coarse-grid variant), mode resampling inflates footprint
and HRR, causing artificially early spread; the documented response gives a
baseline-to-coarse load ratio near 1.86~\cite{Qin2025}.

The integrated-structure treatment groups all cells belonging to one structure,
\[
\bar N=\frac{\sum_{j\in s}N_j}{A_s},\qquad
\mathrm{HRR}_s=q''A_s,
\]
then redistributes load uniformly and emits from the structure centroid.
The integrated-aligned variant expects an aligned \SI{2.5}{m} mesh to
reproduce the baseline. The integrated-misaligned variant permits up to 33\%
load deviation and 10\% ignition-time deviation
on a misaligned \SI{3.24}{m} mesh.

\section{Derivation of expected behavior}
The exact grid should reproduce the analytical structure-load and ignition chronology. Native fine cells should expose within-structure imbalance, native coarse cells should expose footprint inflation, and integrated variants should recover the physical-structure response independently of alignment.

\section{Verification metrics and acceptance criteria}

The postprocessor calculates every metric explicitly. The exact-grid baseline
compares structure-integrated load with the analytical convolution defined
above and second-structure ignition with \SI{2858.8}{s}. The native fine and
coarse variants verify the documented cell-local artifacts rather than claiming
convergence. The integrated aligned and misaligned variants compare
physical-structure profiles with the exact-grid baseline. Missing output or
capability is reported as
\emph{NOT EVALUATED}, never as an implicit pass.


\subsection{Justification of metric quantification}

The criteria intentionally separate reproduction of native-grid defects from
accuracy requirements for the baseline and integrated treatments.  The 10\%
baseline load and ignition bands accommodate rasterized deposition and
threshold-crossing time while remaining small relative to the resolution
effects documented for this configuration~\cite{Qin2025}.  A native fine-grid
load range of at least 5\% is set above roundoff and output quantization, and a
ROS ratio no greater than 0.50 requires the material factor-of-two slowdown
that identifies independent-cell behavior rather than a visually detectable
but negligible difference.  The coarse-grid load ratio is centered on the
reported 1.86 response; its 25\% band permits the severe footprint snapping
intrinsic to a \SI{30}{m} cell without accepting a reversal of the expected
inflation.

For the aligned integrated grid, exact structure boundaries and conservation
of integrated load justify the stringent 1\% load limit.  Its 5\% ignition
limit is wider because ignition is a nonlinear threshold event recorded at a
discrete time.  For the \SI{3.24}{m} misaligned grid, one boundary-cell width
relative to a \SI{10}{m} structure is \(3.24/10=0.324\); the 0.33 maximum
load-error limit therefore admits approximately one unresolved boundary strip
but not a larger structure-scale bias.  The corresponding 10\% ignition bound
allows threshold amplification of that overlap error.  Capability and
completeness are conjunctive because the integrated claims cannot be inferred
from native-grid substitutes.

\subsection{Predeclared acceptance contract}

Table~\ref{tab:acceptance-contract} summarizes the metrics fixed before
inspection of the calculated results. Numerical values and component statuses
are reported separately in Section~7.

\begin{table}[H]
\centering
\footnotesize
\caption{Predeclared verification metrics and acceptance criteria.}
\label{tab:acceptance-contract}
\begin{tabular}{@{}p{0.23\linewidth}p{0.47\linewidth}p{0.21\linewidth}@{}}
\toprule
Metric & Output, region, and aggregation & Acceptance criterion\\
\midrule
Baseline load error & Mean structure-load relative error against the analytical deposition reference, excluding the source structure. & \(\le0.10\) \\
Baseline ignition error & Relative error of second-structure ignition time against \SI{2858.8}{s}. & \(\le0.10\) \\
Native fine-grid load distortion & Median within-structure relative load range. & \(\ge0.05\) \\
Native fine-grid ROS ratio & Fine-grid fitted mean ROS divided by baseline ROS. & \(\le0.50\) \\
Native coarse-grid load ratio & Baseline structure-4 load divided by coarse structure-2 load. & \(1.86\pm25\%\) \\
Integrated aligned load error & Mean matched structure-load relative error against baseline. & \(\le0.01\) \\
Integrated aligned ignition error & Mean matched ignition-time relative error against baseline. & \(\le0.05\) \\
Integrated misaligned load error & Maximum matched structure-load relative error against baseline. & \(\le0.33\) \\
Integrated misaligned ignition error & Maximum matched ignition-time relative error against baseline. & \(\le0.10\) \\
Overall decision & All configured variants and metrics must be evaluable; unavailable integrated capability prevents a PASS. & PASS only if all rows pass \\
\bottomrule
\end{tabular}
\end{table}

\section{Simulation configuration and predicted outcome}

Figure~\ref{fig:input-configuration} records the prepared spatial inputs for \texttt{baseline\_dx10}.
The postprocessor reads these panels from the variant-local GeoTIFF files, removes the
two-cell numerical buffer, and retains map coordinates in metres.  The left panel shows
the most informative fuel or structure layout available for the case; the right panel
shows the cells selected by the non-positive initial level-set field, making a localized
ignition or firebrand-emission source explicit.

\begin{figure}[H]
\centering
\includegraphics[width=0.98\textwidth,height=0.42\textheight,keepaspectratio]{../figures/input_configuration.pdf}
\caption{Whole-domain prepared input configuration for the representative variant.  The
physical domain is shown without the numerical halo; coordinates, configured wind values,
fuel or structure layout, and initial ignition/source geometry are identified in the panels.}
\label{fig:input-configuration}
\end{figure}


\begin{table}[H]
\centering\small
\caption{Five variants generated from the common physical problem.}
\begin{tabular}{llllp{0.31\linewidth}}
\toprule
Variant & \(\Delta x\) (m) & \(t_{\max}\) (s) & Treatment & Expected result\\
\midrule
Baseline & 10.00 & 22000 & native & analytical deposition and ignition\\
Native fine & 2.50 & 50000 & native fine & imbalance and slower ROS\\
Native coarse & 30.00 & 22000 & native coarse & inflated HRR and earlier spread\\
Integrated aligned & 2.50 & 22000 & integrated aligned & baseline-equivalent profile\\
Integrated misaligned & 3.24 & 22000 & integrated misaligned & bounded profile error\\
\bottomrule
\end{tabular}
\end{table}

Each time step is \(\Delta t=0.5\Delta x/U\). All structure rasters reference
a case-local model 14 containing the design fire and physical delay. Structures
form a periodic \(220\times100\)-m two-dimensional lattice. Each represented
cell uses its true square area, so HRRPUA remains \SI{400}{kW.m^{-2}} and the
dimensionless ignition-hardening factor remains one at every resolution. Every
structure in the first upwind column is initialized at \(t=0\).

\subsection{Preprocessing, execution, and postprocessing}

The preprocessor uses exact cell--footprint overlap, applies mode classification
for native fine/coarse experiments, writes structure-ID and overlap rasters,
creates all configurations, and probes the source for integrated capability.
The runner executes supported variants using the absolute executable path in
\texttt{ELMFIRE\_BIN}. The postprocessor reads final ember-flux and TOA
GeoTIFFs, normalizes counts by physical structure area, aggregates by structure
ID, evaluates the formulas above, and writes \texttt{outputs/metrics.json}.

\clearpage
\section{Actual simulation results}

Figure~\ref{fig:domain-result} shows the representative time-of-arrival field from \texttt{baseline\_dx10}.
It is read from the actual ELMFIRE output named in the panel, masks nodata, and excludes
the two-cell numerical buffer.  This completed baseline provides whole-domain evidence even though variants requiring the unavailable integrated-structure capability remain not evaluable.

\begin{figure}[H]
\centering
\includegraphics[width=0.96\textwidth,height=0.50\textheight,keepaspectratio]{../figures/domain_result.pdf}
\caption{Whole-domain time-of-arrival field from the representative completed ELMFIRE run.  This
domain-scale result supplements, rather than replaces, the higher-level verification
profiles, convergence plots, ensemble statistics, and calculated acceptance metrics below.}
\label{fig:domain-result}
\end{figure}


\CompletedVariantCount{} of \RequestedVariantCount{} variants currently have
complete, configuration-matched ELMFIRE output. Source-aligned physical-state
figures are provided for all three native representations; a figure explicitly
states when its current matching output is unavailable. The two
integrated-structure variants remain capability gated.

\begin{figure}[H]
\centering
\includegraphics[width=\linewidth,height=0.68\textheight,keepaspectratio]{baseline_reference.pdf}
\caption{Exact-grid baseline observables: structure load, ignition chronology,
leading-edge trajectory, and instantaneous ROS. Analytical load and
second-structure timing references are overlaid where available.}
\label{fig:wsr-baseline}
\end{figure}

\begin{figure}[H]
\centering
\includegraphics[width=\linewidth,height=0.68\textheight,keepaspectratio]{native_fine_reference.pdf}
\caption{Native fine-grid observables. Cell-scale load variation within each
physical structure is retained so discretization distortion is visible rather
than hidden by structure averaging.}
\label{fig:wsr-fine}
\end{figure}

\begin{figure}[H]
\centering
\includegraphics[width=\linewidth,height=0.68\textheight,keepaspectratio]{native_coarse_reference.pdf}
\caption{Native coarse-grid observables. Mode resampling changes the represented
footprint and separation, and the resulting load and ignition histories are
shown on their actual physical coordinates.}
\label{fig:wsr-coarse}
\end{figure}

Figure~\ref{fig:wsr-summary} is the supplemental cross-variant diagnostic. The
native-test decision is \textbf{\NativeStatus}; the complete-suite decision is
\textbf{\OverallStatus}.

\begin{figure}[H]
\centering
\includegraphics[width=\linewidth]{wui_structure_resolution.pdf}
\caption{Structural-resolution summary. (a) Rasterized layouts; (b) deposited
load by physical structure; (c) ignition chronology; and (d) component decision
status, with unsupported components shown as not evaluated.}
\label{fig:wsr-summary}
\end{figure}

\section{Calculated metrics and verification decision}

\begin{table}[H]
\centering\small
\caption{Calculated verification metrics and component decisions.}
\begin{tabular}{@{}p{0.37\linewidth}p{0.18\linewidth}p{0.17\linewidth}p{0.14\linewidth}@{}}
\toprule
Metric & Acceptance limit & Calculated & Status\\
\midrule
\MetricRows
\bottomrule
\end{tabular}
\end{table}

The cited model description~\cite{Qin2025} states that the integrated-structure method
had not yet been implemented in ELMFIRE. The preprocessor confirms this against
the current source. The integrated aligned and misaligned cases are therefore
generated as prospective variants
but are not run while that feature is absent. This prevents the former
resolution-independent directories from silently rerunning the native
cell-local algorithm and reporting misleading results. Pre-merge input snapshots are archived outside the canonical case directory and are excluded from the run plan.

\section{Technical interpretation}
The completed exact-grid baseline fails both of its quantitative reference
checks. The native fine-grid run is long and has no completed matching output,
while the retained coarse output is rejected because its filename encodes a
different terminal time than the current manifest. Those components are
therefore not evaluated. The integrated variants test whether one physical
structure remains the source and collector; their missing capability is a
source-interface limitation and is likewise reported as not evaluable rather
than as a numerical pass or failure.

\section{Scope and limitations}
The case is an idealized two-dimensional WUI lattice with fixed structure geometry and wind. A representative cross-stream lattice row supplies the one-dimensional structure profiles, while whole-domain maps expose unintended lateral behavior. The case verifies resolution behavior and capability coverage, not the physical accuracy of the structural transport or ignition correlations.

\begin{thebibliography}{9}
\bibitem{Qin2025}
Y. Qin, \emph{A Physics-Based Eulerian Framework for Modeling Firebrand
Showering in Regional-Scale Wildland and Wildland--Urban Interface Fire
Simulations}, Ph.D. dissertation, University of Maryland, College Park, 2025.
\end{thebibliography}
